% Figure rendered 5.5in wide with 8pt labels, i.e. \linewidth in the ICLR template.
% Same held-out slices as 00_inpaper_fig_compare (fig:adni-compare), which follows it.
\begin{figure}[t]
\centering
\includegraphics[width=\linewidth]{adni/00_inpaper_fig_ablation.pdf}
\caption{Ablation of MMCLAST-cg on ADNI. Each row is a held-out test slice from
a different subject; the top three rows translate T1 to FA and the bottom three
FA to T1. Columns show the source image, the target, and the variants of
Table~\ref{tab:recon}: \emph{base} has neither $\mathcal{L}_{\text{latcyc}}$ nor
$\mathcal{L}_{\text{path}}$; \emph{latcyc} adds $\mathcal{L}_{\text{latcyc}}$;
\emph{morph} further adds $\mathcal{L}_{\text{path}}$ (adversarial term with
$D_{\text{mix}}$ plus smoothness term); \emph{morph-smooth} keeps only the
smoothness term; \emph{morph-bi} applies $\mathcal{L}_{\text{path}}$ in both
directions. All variants recover the same anatomy, and the terms that supervise
the path change the endpoint only in texture: latcyc leaves more high-frequency
speckle in the synthesized FA maps, while morph-bi produces the smoothest
tracts, in line with its highest T1$\rightarrow$FA SSIM in
Table~\ref{tab:recon}. The FA$\rightarrow$T1 outputs are nearly
indistinguishable across variants.}
\label{fig:adni-ablation}
\end{figure}
